% Formal proof-method section; these generic and source-certificate theorems
% were compiled in the isolated K227F and K226F namespaces.
\section{A checked source equivalence with one register offset}
\label{sec:offset-bisimulation}

The increment-and-drain rotation changes the value of register 5 at an
intermediate test. At the corresponding \texttt{CD6} checkpoints, the new
program has already performed its extra increment, while the old program
performs that increment after leaving the loop. Equality of register files
at every test is therefore an unnecessarily strong certificate format.
The source bridge permits a single, explicit integer offset.

An entry of the finite certificate has the form
\[
  (P_{\mathrm{n}},P_{\mathrm{o}},z,j_+,j_0),
\]
where \(P_{\mathrm{n}},P_{\mathrm{o}}\) are the new and old instruction-pointer bit
lists, \(z\in\mathbb Z\), and \(j_+,j_0\) identify the entries reached
after positive and zero tests. The invariant at that entry is
\begin{equation}
  R_{\mathrm{n}}(i)=R_{\mathrm{o}}(i)+\delta_z(i),
  \qquad
  \delta_z(i)=\begin{cases}z&i=5,\\0&i\ne5.\end{cases}
  \label{eq:register-offset}
\end{equation}
The equality is interpreted in \(\mathbb Z\); both register files
themselves remain functions to \(\mathbb N\). No negative register value
is introduced into either program.

Between conditional decrements, the inherited \texttt{seg} function runs
increments and control instructions until it reaches another conditional
decrement or a halt. Let \(c_{\mathrm{n}}(i),c_{\mathrm{o}}(i)\) be the numbers of
increments of register \(i\) in the two segments. A branch leading to an
entry with offset \(z'\) passes precisely when the endpoint pointers match
and, for every register that occurs in either segment and for register 5,
\begin{equation}
  \delta_z(i)+c_{\mathrm{n}}(i)
       =c_{\mathrm{o}}(i)+\delta_{z'}(i).
  \label{eq:offset-counts}
\end{equation}
Registers absent from this finite list have zero counts and zero relevant
offset, so the check covers all register indices.

Both checkpoints must test the same register \(r\), and the checker
requires \(\delta_z(r)=0\). Equation~\eqref{eq:register-offset} therefore
gives equal tested values, equal test outcomes, and equal decremented values
on the positive branch. Equation~\eqref{eq:offset-counts} then establishes
the next checkpoint's invariant. A halt is matched to a halt. The initial
segments start with offset zero.

\texttt{OffsetBisim230.v} proves soundness by strong induction on the
number of steps in a finite halting execution. It reuses the pinned
\texttt{Bisim230.v} proofs of segment execution, and supplies the additional
integer arithmetic and zero-offset test argument. Thus it does not assume
that the programs terminate. The mirror certificate swaps the pointers
and negates every offset; checking it gives the converse implication.
The exported theorem \texttt{ocheck\_iff} proves halting equivalence for
every initial natural-valued register file.

For the 227-state source, the checked certificate contains 2373 entries
and 4744 conditional branches. For the 226-state source, it contains
2376 entries and 4750 branches. Each has only eight entries with offset 1:
two instruction-pointer spellings for each of the four rotated loop
occurrences. Every nonzero-offset entry tests register 6. Every segment
has length at most five. Both directions of both certificates evaluate
to true in Rocq, yielding closed source-equivalence theorems
\texttt{p227\_p230} and \texttt{p226\_p230}. These apply to the actual
bootstrap register file as a special case.

\section{Literal runtime replay and exact state counting}
\label{sec:literal-runtime-replay}

The runtime proof can be checked directly against a modified table. The
isolated proof namespaces retain K230's sparse state labels, disable every
deleted label by giving it no transition, and replay the literal bootstrap,
the four register-block proofs, every interpreter leaf, and the global
simulation invariant. Original source, logic, and set-theory proof files
are reused without modification. The new source-equivalence theorem then
connects the changed interpreter tree to the pinned \(p_{230}\) source.
This route discharges native context requirements inside the runtime proof,
without assuming a separately stated global guard theorem.

Sparse labels serve only to keep inherited local tape proofs readable.
The delivered machine uses dense labels. Its separate Rocq context has
state type \(\{q\in\mathbb N:q\le N\}\), alphabet \(\{0,1\}\),
working states \(0,\ldots,N-1\), and terminal state \(N\).
The dense literal array has exactly \(N\) paired rows (two transitions each), and a finite checker
verifies every target is at most \(N\).

Let \(e\) be the increasing map that skips deleted sparse labels.
A finite proof checks the transition identity for every dense state and
both symbols, including the terminal state: applying the sparse machine
at state \(e(q)\) is exactly the dense transition with its target state
mapped through \(e\). Tape symbols and head movements agree. The induced
map on configurations changes only the state. Induction transports every
dense execution to a sparse execution, while the deterministic transition
identity and the positive-progress halting lemma give the reverse halting
implication. The transition argument treats arbitrary tapes at corresponding states.
The exported halting equivalence is its blank instance, which composes
with the runtime and source proofs.

Consequently the state count is a property of the exact delivered table
and of its own finite Rocq context. It is not a count obtained by leaving
an unexecuted working state in the formal machine and omitting it from
the displayed census.


The exact endpoint names, assumption reports, and successful independent
kernel recheck are recorded in Section~\ref{sec:verification}.
